\documentclass[11pt]{article}

\newtheorem{theorem}{Theorem}
\newtheorem{lemma}{Lemma}
\newenvironment{proof}{\par\noindent{\it Proof. }}{\hfill$\diamond$\par}

\newcommand{\Applicants}{\mathcal A}
\newcommand{\Slots}{\mathcal S}
\newcommand{\chosen}{C}
\newcommand{\border}{V_{\chosen}}
\newcommand{\weight}{w}
\newcommand{\class}{\kappa}

\title{Clean Proof of the LAP Variability Bound}
\author{Lean formalization note for DGD26}
\date{}

\begin{document}
\maketitle

\paragraph{Purpose.}
This note records the paper-facing proof of the finite linear-assignment
variability result.  The source paper states this result at a high level.  The
Lean development supplies the detailed finite exchange argument: exact
one-for-one insertion, directed alternating splices, a tail-dominance inequality,
same-order injectivity for borderline applicants, and the final counting bound.

\section*{Model and theorem}

Let $\Applicants$ be a finite applicant universe and let $\Slots$ be a finite
set of assignment slots.  A feasible assignment for an applicant pool
$X \subseteq \Applicants$ assigns each slot either no applicant or one applicant
from $X$, and assigns no applicant to two slots.  In the capacity-filling regime
used here, a feasible assignment fills all slots when $|X| \ge |\Slots|$.

Each applicant-slot pair has an integer weight $\weight(a,s)$.  A selector
chooses, for each pool $X$, a unique global optimum feasible capacity-filling
assignment $A_X$.  The induced choice rule is
\[
  \chosen(X) = \{a : a \hbox{ is assigned in } A_X\}.
\]
For a fixed slot $s$, write $a <_s b$ when
\[
  \weight(a,s) < \weight(b,s),
\]
so $b$ is strictly preferred to $a$ at slot $s$.  Assume every slot has no ties:
for distinct applicants $a,b$, exactly one of $a <_s b$ or $b <_s a$ holds.

Let $\class : \Slots \to K$ be any classifier such that
\[
  \class(s)=\class(t)
  \quad\Longrightarrow\quad
  (a <_s b \iff a <_t b)
\]
for every pair of applicants $a,b$.  Thus each classifier value is a slot-order
class.

\begin{theorem}[linear-assignment variability bound]
For every pool $X$,
\[
  |\border(X)| \le |\class(\Slots)|,
\]
where $\border(X)$ is the set of applicants in $\chosen(X)$ who can be displaced
by adding one fresh applicant.  Equivalently, the assignment-induced choice rule
has variability bounded by the number of represented slot-order classes.
\end{theorem}

\section*{Exact exchange and directed paths}

Fix a pool $X$, a fresh applicant $x \notin X$, and let
\[
  A = A_X,\qquad B = A_{X\cup\{x\}}.
\]
Suppose $y \in \chosen(X)\setminus \chosen(X\cup\{x\})$.  The formalization first
proves the exact one-for-one exchange property: under the unique-optimum and
capacity-filling hypotheses, the new chosen set is
\[
  \chosen(X\cup\{x\}) = (\chosen(X)\setminus \{y\})\cup\{x\}.
\]
In particular every old chosen applicant other than $y$ remains chosen after
the insertion.

The old and new assignments define a directed alternating slot graph: there is
an edge from slot $s$ to slot $t$ if the old occupant of $s$ in $A$ is assigned
to $t$ in $B$.  Let $s_0$ be the new slot containing $x$ in $B$ and let $s_k$ be
the old slot containing $y$ in $A$.  The one-for-one exchange yields a directed
path
\[
  s_0,s_1,\ldots,s_k
\]
from the fresh slot to the lost slot.  Write
\[
  a_i = A(s_i) \qquad (0\le i\le k),
\]
so $a_k=y$, and the path identities are
\[
  B(s_0)=x,\qquad B(s_{i+1})=a_i \quad (0\le i<k).
\]

\section*{Two exchange lemmas}

\begin{lemma}[dominated old slots lie on the fresh-to-lost path]
Let $s$ be an old slot with old occupant $z=A(s)$.  If
\[
  z <_s y,
\]
then $s$ is on the directed path from $s_0$ to $s_k$.
\end{lemma}

\begin{proof}
Suppose $s$ is not on that path.  Splice the new assignment $B$ into the old
assignment $A$ on the fresh-root reachable component and keep $A$ elsewhere.
The finite component-closure facts verified in Lean show that the splice is a
feasible old-pool assignment, still capacity-filling, and still globally
optimal: the complementary splice is feasible for the enlarged pool, and the
two objective inequalities cancel.

Because $s$ is outside the fresh-root component, the spliced old-pool optimum
still assigns $z$ to $s$.  Because $y$ is the unique old applicant lost in
$B$, the same splice still rejects $y$.  Replacing $z$ by $y$ at slot $s$ would
be a feasible one-slot change and would strictly increase the objective, since
$z<_s y$.  This contradicts local optimality, which is derived in Lean from
global objective optimality.  Therefore $s$ must lie on the fresh-to-lost path.
\end{proof}

\begin{lemma}[tail dominance along the path]
For every $i<k$,
\[
  \weight(y,s_i) \le \weight(a_i,s_i).
\]
\end{lemma}

\begin{proof}
Fix $i<k$.  Define
\[
  P_i=\sum_{j=i}^{k-1}\weight(a_j,s_{j+1}),
  \qquad
  Q_i=\sum_{j=i+1}^{k}\weight(a_j,s_j).
\]
Construct an old-pool competitor $E_i$ from $A$ by rotating the suffix
$s_i,s_{i+1},\ldots,s_k$ forward and putting the lost applicant $y=a_k$ into
$s_i$:
\[
  E_i(s_i)=y,\qquad
  E_i(s_{j+1})=a_j \quad (i\le j<k),
\]
and $E_i$ agrees with $A$ outside the suffix.  This is feasible because it only
permutes applicants already used by $A$ on distinct suffix slots.  Since $A$ is
old-pool optimal and the outside terms cancel,
\[
  \weight(y,s_i)+P_i \le \weight(a_i,s_i)+Q_i. \qquad (1)
\]

Construct an enlarged-pool competitor $F_i$ from $B$ by setting the tail
$s_{i+1},\ldots,s_k$ back to the old assignment:
\[
  F_i(s_j)=a_j \quad (i<j\le k),
\]
and leaving $B$ unchanged outside that tail.  This drops $a_i$ from the shifted
tail and inserts $y$ at the lost slot, so feasibility follows from the exact
one-for-one exchange and no-duplicate assignment facts.  Since $B$ is
enlarged-pool optimal and outside terms cancel,
\[
  Q_i \le P_i. \qquad (2)
\]

Combining (1) and (2) gives
\[
  \weight(y,s_i)+P_i
  \le
  \weight(a_i,s_i)+Q_i
  \le
  \weight(a_i,s_i)+P_i.
\]
Canceling $P_i$ proves $\weight(y,s_i)\le \weight(a_i,s_i)$.
\end{proof}

\section*{Same-order injectivity}

\begin{lemma}[same-order borderline injectivity]
Fix $X$.  Let $y,z\in \border(X)$, and suppose the selected old assignment
$A_X$ assigns $y$ to slot $s_y$ and $z$ to slot $s_z$.  If
$\class(s_y)=\class(s_z)$, then $y=z$.
\end{lemma}

\begin{proof}
Assume $y\ne z$.  By no ties at slot $s_z$, either $z<_{s_z}y$ or
$y<_{s_z}z$.

If $z<_{s_z}y$, use the fresh insertion witnessing $y\in\border(X)$, so $y$ is
the unique old applicant lost.  Since $z$ remains chosen and is below $y$ at
its old slot $s_z$, the dominated-slot lemma puts $s_z$ on the fresh-to-lost
path.  Applying tail dominance at that path position yields
\[
  \weight(y,s_z)\le \weight(z,s_z),
\]
contradicting $z<_{s_z}y$.

If $y<_{s_z}z$, the common slot-order hypothesis transports the same strict
comparison to $s_y$, so $y<_{s_y}z$.  Now use the fresh insertion witnessing
$z\in\border(X)$, with $z$ as the unique old applicant lost.  The same argument
puts $s_y$ on the corresponding fresh-to-lost path and gives
\[
  \weight(z,s_y)\le \weight(y,s_y),
\]
contradicting $y<_{s_y}z$.

Both cases are impossible, so $y=z$.
\end{proof}

\section*{Counting}

For each borderline applicant $y\in\border(X)$, let $s_y$ be the unique slot
where $A_X$ assigns $y$.  Define
\[
  \phi(y)=\class(s_y).
\]
The map is well-defined because $y\in\chosen(X)$ and assignments do not
duplicate applicants.  The same-order injectivity lemma says exactly that
$\phi$ is injective.  Since both sets are finite,
\[
  |\border(X)| \le |\class(\Slots)|.
\]
This proves the linear-assignment variability bound.

\end{document}
